\scanpage{18}
\mspage{9}
\uline{Définition de $\alpha$}\qquad On va définir
\[
\alpha(\mathcal E,\mathcal G,\pi,Z,\ldots)=(\Pi,\widetilde\Pi,p,\mathcal G,a)
\]
\[
\left\{\begin{array}{l}
\operatorname{Ob}\Pi=Z/\pi\quad(\text{soit }J)\\
\operatorname{Isom}(j,j')=\operatorname{Hom}_{\pi}(Z_j,Z_{j'})\\
\text{composition des flèches évidente}
\end{array}\right.\qquad\bigg|\quad\text{d'où groupoïde }\Pi
\]
\[
\deleted{\widetilde\Pi_j=Z_j}\qquad\operatorname{Ob}\widetilde\Pi=Z
\]
\diagram{\begin{tikzcd}[ampersand replacement=\&]
\operatorname{Isom}(z,z')=\{\unsure{k|*}\}\arrow[d]\&\begin{gathered}\forall z,z'\in Z\\(\widetilde\Pi\text{ cat. discrète})\end{gathered}\\
\operatorname{Isom}(j,j')\simeq\operatorname{Hom}(Z_j,Z_{j'})\&
\end{tikzcd}}
\deleted{(\unsure{où} $j,j'$ \unsure{les} \illegible{1})}\inserted{(où $j=\dot z$, $j'=\dot z'$)}
\margin{$\Longleftarrow$ $\widetilde\Pi$ cat. \uline{discrète} sur $\Pi$ définie par foncteur \inserted{contr.}
\diagram{\begin{tikzcd}[ampersand replacement=\&]
F:\Pi^\circ \arrow[r] \& \operatorname{Ens}
\end{tikzcd}}
$F(j)=Z_j,\ldots$\quad $\widetilde\Pi\simeq\Pi/F$}
$\downarrow$ transforme l'unique pt de $\operatorname{Hom}(z,z')$ en l'unique \inserted{$\pi$-}iso $Z_j\mathbin{-}Z_{j'}$ transformant $z$ en $z'\ldots.$

Pour l'instant on n'a pas utilisé $\mathcal E$ et l'op. de $\mathcal E$ sur $Z$. On va l'utiliser maintenant pour faire opérer $\mathcal G=\mathcal E/\pi$ sur $\Pi$, et $\mathcal E$ sur $\widetilde\Pi$, de façon compatible.
Comme $\mathcal E$ opère sur $Z=\operatorname{Ens}(\widetilde\Pi)$, il opère \unsure{en fait} sur $\widetilde\Pi$, cat. discrète associée, \deleted{\unsure{opérant librement}} et \deleted{le} $\pi$ opère par $\Pi$-auto---, d'où l'opération de $\mathcal G$ sur $\unsure{K|\Pi}$ par passage au quotient.
\mspage{17}
